\chapter{Αλγόριθμοι Επίλυσης}
\label{chap:algorithms}

Το κεφάλαιο αυτό παρουσιάζει τη σχεδίαση των αλγορίθμων επίλυσης της παρούσας εργασίας πάνω στο μοντέλο κατάστασης του Κεφαλαίου~\ref{chap:model}. Εξετάζονται οι κοινές αρχές όλων των επιλυτών, η BFS ως σημείο αναφοράς, η αλυσίδα υποομάδων του Thistlethwaite, ο αλγόριθμος δύο φάσεων του Kociemba, η βέλτιστη αναζήτηση κατά Korf με IDA*, η πλήρης μελέτη του κύβου 2×2×2 και η εντολή αναγνώρισης απόστασης. Η αντίστοιχη υλοποίηση περιγράφεται στο Κεφάλαιο~\ref{chap:implementation} και η πειραματική αξιολόγηση στο Κεφάλαιο~\ref{chap:experiments}.

\section{Κοινές αρχές σχεδίασης}
\label{sec:algorithms-principles}
Οι αλγόριθμοι της εργασίας δεν υλοποιούνται ως ανεξάρτητα προγράμματα με διαφορετικές συμβάσεις. Όλοι οι επιλυτές (\eng{solvers}) δέχονται την ίδια κατάσταση κυβιδίων (\eng{cubies}), χρησιμοποιούν τον ίδιο αναλυτή κινήσεων στη μετρική μισών στροφών (\eng{half-turn metric}, HTM), επιστρέφουν τυποποιημένο αποτέλεσμα και περνούν από τον ίδιο ανεξάρτητο επαληθευτή (\eng{verifier}). Η κοινή αυτή υποδομή είναι απαραίτητη για δίκαιη σύγκριση. Αν ένας επιλυτής μετρούσε το \(R2\) ως δύο κινήσεις και άλλος ως μία, ή αν ένας επιλυτής δεχόταν μη φυσικές καταστάσεις, οι πίνακες αποτελεσμάτων δεν θα ήταν συγκρίσιμοι.

Κάθε επιλυτής επιστρέφει χαρακτηρισμό κατάστασης (\eng{status}) και όχι μόνο ακολουθία κινήσεων. Η σχεδίαση αυτή είναι συνειδητή. Στην αλγοριθμική επίλυση του κύβου υπάρχουν διαφορετικά επίπεδα εγγύησης: πλήρης απόδειξη ελάχιστης απόστασης (\code{exact}), επαληθευμένη λύση χωρίς απόδειξη ελαχιστότητας (\code{non_exact}), αποδεδειγμένο κάτω φράγμα (\code{lower_bound}), εξάντληση των ορίων αναζήτησης (\code{timeout}) και μη εφαρμόσιμη μέθοδος. Αν όλες αυτές οι περιπτώσεις συμπτυχθούν σε ένα απλό πεδίο επιτυχίας, η εργασία χάνει ακριβώς τη διάκριση που ζητά το θέμα. Γι' αυτό η αρχιτεκτονική των επιλυτών κάνει την αβεβαιότητα ορατή.

Η δεύτερη κοινή αρχή είναι ότι τα κάτω φράγματα πρέπει να είναι παραδεκτά όταν χρησιμοποιούνται για ακριβή αναζήτηση. Η IDA* μπορεί να χαρακτηρίσει αποτέλεσμα ως \code{exact} μόνο όταν η ευρετική δεν υπερεκτιμά και η αναζήτηση ολοκληρώνει το αντίστοιχο όριο. Για πρακτικούς επιλυτές διαδοχικών σταδίων, όπως οι Kociemba και Thistlethwaite, μια λύση μπορεί να είναι απολύτως έγκυρη και χρήσιμη χωρίς να είναι αποδεδειγμένα ελάχιστη. Η εργασία επομένως επιτρέπει το \code{non_exact} ως θετικό, επαληθευμένο αποτέλεσμα, αλλά δεν το ονομάζει βέλτιστο.

\section{Η BFS ως σημείο αναφοράς}
\label{sec:algorithms-bfs}
Η BFS χρησιμοποιείται για μικρά βάθη επειδή έχει απλή και ισχυρή ιδιότητα: αν εξετάσει όλα τα επίπεδα μέχρι το βάθος \(d-1\) και βρει πρώτη λύση στο βάθος \(d\), τότε η απόσταση είναι ακριβώς \(d\). Στην παρούσα εργασία η BFS δεν αποτελεί γενικό επιλυτή για τον κύβο 3×3×3. Λειτουργεί ως σημείο αναφοράς για τις ρηχές περιπτώσεις του σώματος μετρήσεων (\eng{benchmark corpus}) και ως μηχανισμός παραγωγής της πλήρους κατανομής αποστάσεων του κύβου 2×2×2 (\eng{Pocket Cube}).

Η σχεδίαση της BFS αποφεύγει πλεονασμούς με τον ίδιο κανόνα που χρησιμοποιείται και στην IDA*: δεν παράγονται διαδοχικές κινήσεις στην ίδια έδρα. Η αποφυγή αυτή δεν αλλάζει την ορθότητα για την αναζήτηση μικρότερης ακολουθίας, επειδή κάθε ζεύγος διαδοχικών κινήσεων στην ίδια έδρα μπορεί να αντικατασταθεί από μία ισοδύναμη κίνηση ή να απαλειφθεί. Μειώνει όμως τον αριθμό των παραγόμενων κόμβων και καθιστά τις ρηχές μετρήσεις πιο σταθερές.

Στα αποτελέσματα, η BFS εμφανίζεται έμμεσα ως πεδίο \code{known_exact_distance} (γνωστή ακριβής απόσταση) για περιπτώσεις όπου το βάθος βρίσκεται μέσα στο όριο. Η τιμή αυτή χρησιμοποιείται και από τον επαληθευτή αποτελεσμάτων: αν ένας επιλυτής δηλώσει \code{exact} και η BFS έχει υπολογίσει ακριβή ρηχή απόσταση, το μήκος της λύσης πρέπει να συμφωνεί. Έτσι η BFS λειτουργεί ως ανεξάρτητος έλεγχος και όχι μόνο ως ένας ακόμη επιλυτής.

\section{Ο αλγόριθμος Thistlethwaite}
\label{sec:algorithms-thistlethwaite}
Η μέθοδος Thistlethwaite μειώνει την κατάσταση μέσα από αλυσίδα υποομάδων. Η τεκμηρίωση του Kociemba περιγράφει ότι η μέθοδος χρησιμοποιεί ενδιάμεσες υποομάδες και στατικούς πίνακες ελιγμών, με άνω φράγμα 52 κινήσεων \cite{kociembaImplementationDetails}. Δευτερογενής τεχνική περιγραφή της αλυσίδας υποομάδων και των φάσεων δίνεται επίσης από τον Scherphuis \cite{scherphuisThistlethwaite}. Η παρούσα εργασία υλοποιεί εγγενή τετραφασική εκδοχή της ιδέας: \(G_0 \rightarrow G_1\) για τον προσανατολισμό ακμών, \(G_1 \rightarrow G_2\) για τους προσανατολισμούς και τη φέτα UD (\eng{UD-slice}), \(G_2 \rightarrow G_3\) για την υποομάδα των μισών στροφών (\eng{square group}) και \(G_3 \rightarrow I\) για την τελική λύση με μισές στροφές. Η υλοποίηση δεν παρουσιάζεται ως αναπαραγωγή των ιστορικών στατικών πινάκων ελιγμών του Thistlethwaite· στη θέση τους χρησιμοποιεί ακριβείς πίνακες αποστάσεων BFS ανά φάση. Οι πίνακες αυτοί περιέχουν την πραγματική απόσταση κάθε φασικής κατάστασης από την επόμενη υποομάδα. Έτσι η μέθοδος τερματίζει εγγυημένα σε κάθε έγκυρη κατάσταση, χωρίς όρια χρόνου ή κόμβων. Κάθε λύση που περνά την ανεξάρτητη επαλήθευση σημειώνεται ως \code{non_exact}. Ο μόνος άλλος δυνατός χαρακτηρισμός είναι \code{failed}, αν η επαλήθευση αποτύχει· αποτελέσματα \code{timeout} ή \code{lower_bound} δεν εμφανίζονται σε αυτόν τον επιλυτή.

Κάθε μείωση υποομάδας καθοδηγείται από προϋπολογισμένο, ακριβή πίνακα αποστάσεων BFS πάνω στον αντίστοιχο χώρο συντεταγμένων (\eng{coordinates}) ή αριστερών συμπλόκων (\eng{cosets}). Για τη φάση \(G_0 \rightarrow G_1\) ο πίνακας ορίζεται στη συντεταγμένη προσανατολισμού ακμών (\(2^{11}=2048\) καταστάσεις) με το πλήρες σύνολο των 18 κινήσεων HTM, επειδή δεν υπάρχουν ακόμη περιορισμοί που πρέπει να διατηρηθούν. Για τη φάση \(G_1 \rightarrow G_2\) ο πίνακας ορίζεται στο γινόμενο του προσανατολισμού γωνιών (\(3^7=2187\)) και του συνδυασμού της φέτας UD (\(\binom{12}{4}=495\)), συνολικά \(1\,082\,565\) καταστάσεις, με το σύνολο κινήσεων που διατηρεί τον προσανατολισμό ακμών: \(U,D,R,L\) σε όλες τις παραλλαγές HTM και \(F2,B2\). Για τη φάση \(G_2 \rightarrow G_3\) ο πίνακας ορίζεται στο γινόμενο των αριστερών συμπλόκων της μετάθεσης γωνιών ως προς τη δράση της υποομάδας μισών στροφών (420 σύμπλοκα) και των αντίστοιχων συμπλόκων της μετάθεσης ακμών (140 σύμπλοκα), με τα δέκα φασικά σύμβολα HTM \(U,U',U2,D,D',D2,R2,L2,F2,B2\). Από τους \(58\,800\) συνδυασμούς δεικτών, οι \(29\,400\) αντιστοιχούν σε προσβάσιμα σύμπλοκα. Για τη φάση \(G_3 \rightarrow I\) ο πίνακας καλύπτει ολόκληρη την υποομάδα μισών στροφών (\(96\cdot 6912=663\,552\) καταστάσεις) με τις έξι μισές στροφές. Η αλυσίδα των τεσσάρων σταδίων, τα μεγέθη των πινάκων και τα μέγιστα φασικά βάθη συνοψίζονται στο Σχήμα~\ref{fig:thistlethwaite-chain}.

\begin{figure}[htbp]\centering
% Αλυσίδα υποομάδων Thistlethwaite (περιεχόμενο σχήματος· το περιβάλλον figure
% και η λεζάντα ορίζονται στο κεφάλαιο 4). Μεγέθη πινάκων και μέγιστα βάθη
% κατά τα κλειδωμένα δεδομένα: 2048 / 1 082 565 / 29 400 προσβάσιμες (από
% 58 800) / 663 552 καταχωρίσεις, βάθη <= 7/10/13/15.
\resizebox{\textwidth}{!}{%
\begin{tikzpicture}[
  gbox/.style={draw, rectangle, align=center, font=\footnotesize,
               minimum height=1.35cm, minimum width=1.8cm, inner sep=4pt},
  stage/.style={font=\footnotesize},
  ann/.style={font=\scriptsize, align=center, anchor=north},
  flow/.style={-{Stealth[length=2.6mm]}, semithick},
  node distance=0.95cm
]
\node[gbox] (g0) {\(G_0\)\\ πλήρης ομάδα};
\node[gbox, right=of g0] (g1) {\(G_1\)\\ προσανατολισμένες\\ ακμές};
\node[gbox, right=of g1] (g2) {\(G_2\)\\ προσανατολισμοί\\ και φέτα UD};
\node[gbox, right=of g2] (g3) {\(G_3\)\\ υποομάδα\\ μισών στροφών};
\node[gbox, right=of g3] (g4) {\(G_4 = I\)\\ λυμένη\\ κατάσταση};

\draw[flow] (g0) -- node[stage, above] {Στάδιο 1} (g1);
\draw[flow] (g1) -- node[stage, above] {Στάδιο 2} (g2);
\draw[flow] (g2) -- node[stage, above] {Στάδιο 3} (g3);
\draw[flow] (g3) -- node[stage, above] {Στάδιο 4} (g4);

\node[ann] at ($(g0.east)!0.5!(g1.west) + (0,-0.8)$)
  {2048 καταχωρίσεις\\ βάθος \(\le 7\)};
\node[ann] at ($(g1.east)!0.5!(g2.west) + (0,-0.8)$)
  {1\,082\,565 καταχωρίσεις\\ βάθος \(\le 10\)};
\node[ann] at ($(g2.east)!0.5!(g3.west) + (0,-0.8)$)
  {29\,400 προσβάσιμες\\ (από 58\,800)\\ βάθος \(\le 13\)};
\node[ann] at ($(g3.east)!0.5!(g4.west) + (0,-0.8)$)
  {663\,552 καταχωρίσεις\\ βάθος \(\le 15\)};
\end{tikzpicture}%
}

\caption{Η αλυσίδα υποομάδων του αλγορίθμου Thistlethwaite: από την πλήρη ομάδα \(G_0\) έως τη λυμένη κατάσταση, με το μέγεθος του πίνακα αποστάσεων κάθε σταδίου (2048, 1\,082\,565, 29\,400 προσβάσιμες από 58\,800, 663\,552 καταχωρίσεις) και τα μέγιστα βάθη 7/10/13/15.}
\label{fig:thistlethwaite-chain}
\end{figure}

Επειδή κάθε πίνακας αποθηκεύει την πραγματική απόσταση μέχρι την επόμενη υποομάδα, η επίλυση κάθε φάσης δεν απαιτεί αναζήτηση. Είναι άπληστη κάθοδος καθοδηγούμενη από τον πίνακα: σε κάθε βήμα υπάρχει εγγυημένα κίνηση που μειώνει την αποθηκευμένη απόσταση κατά ένα. Έτσι η φάση ολοκληρώνεται σε ακριβώς τόσες κινήσεις όση η αρχική τιμή του πίνακα, δηλαδή στο ελάχιστο δυνατό μήκος για τη συγκεκριμένη μετάβαση υποομάδας με το επιτρεπτό σύνολο κινήσεων. Τα μέγιστα φασικά βάθη που προκύπτουν από τους πίνακες είναι 7, 10, 13 και 15 κινήσεις και συμπίπτουν με τα δημοσιευμένα φασικά φράγματα της μεθόδου \cite{scherphuisThistlethwaite}, γεγονός που λειτουργεί ως ανεξάρτητος έλεγχος ορθότητας των χώρων συντεταγμένων και συμπλόκων. Το συνολικό μήκος λύσης φράσσεται επομένως από \(7+10+13+15=45\) κινήσεις.

Η εγγύηση αυτή είναι φασική και όχι καθολική. Η συνένωση φασικά βέλτιστων τμημάτων δεν αποτελεί εν γένει βέλτιστη λύση στον πλήρη χώρο καταστάσεων, επειδή κάθε στάδιο περιορίζεται στο δικό του σύνολο κινήσεων και δεν εξετάζει εναλλακτικές ενδιάμεσες καταστάσεις. Γι' αυτό κάθε λύση χαρακτηρίζεται \code{non_exact} ακόμη και όταν επαληθεύεται πλήρως, και η μέθοδος δεν υποκαθιστά τη βέλτιστη αναζήτηση.

Ο επιλυτής κατασκευάζει την τελική λύση με συνένωση των λύσεων των φάσεων. Η συνένωση είναι έγκυρη μόνο αν κάθε φάση εφαρμοστεί πάνω στην κατάσταση που προκύπτει από τις προηγούμενες κινήσεις. Γι' αυτό η υλοποίηση, μετά από κάθε φάση, εφαρμόζει τη λύση στην κατάσταση κυβιδίων και εκτελεί την επόμενη φάση στη νέα κατάσταση. Στο τέλος ο ανεξάρτητος επαληθευτής εφαρμόζει ολόκληρη τη συνενωμένη ακολουθία από την αρχική κατάσταση. Αν οποιαδήποτε φάση είχε λανθασμένη σύμβαση κινήσεων, ο τελικός έλεγχος θα αποτύγχανε.

\section{Ο αλγόριθμος δύο φάσεων του Kociemba}
\label{sec:algorithms-kociemba}
Ο Kociemba αντικαθιστά την αλυσίδα πολλών ενδιάμεσων ομάδων με μία βασική ενδιάμεση ομάδα \(G_1=\langle U,D,R2,L2,F2,B2\rangle\). Η γραφή αυτή είναι γεννητορική: οι γεννήτορες \(U\) και \(D\) παράγουν και τις τρεις HTM παραλλαγές τους, ενώ οι \(R,L,F,B\) επιτρέπονται μόνο ως μισές στροφές. Στην πρώτη φάση ζητείται κατάσταση με μηδενικούς προσανατολισμούς και τις τέσσερις ακμές της φέτας UD στη σωστή φέτα, ενώ στη δεύτερη φάση η λύση ολοκληρώνεται εντός της περιορισμένης ομάδας \cite{kociembaTwoPhase}. Επειδή ένα μεγαλύτερο μήκος φάσης 1 μπορεί να οδηγήσει σε μικρότερο συνολικό μήκος, η μέθοδος εξετάζει με επαναληπτική εμβάθυνση διαδοχικά μήκη φάσης 1 και παράγει διαδοχικά καλύτερες συνολικές λύσεις \cite{kociembaTwoPhase}. Η δομή των δύο φάσεων απεικονίζεται στο Σχήμα~\ref{fig:kociemba-twophase}. Η πρακτική γραμμή του Cube Explorer δείχνει ότι η μέθοδος μπορεί να υλοποιηθεί ως ιδιαίτερα αποδοτικό λογισμικό επίλυσης, αλλά η απόδοση αυτή βασίζεται σε λεπτομερείς πίνακες και βελτιστοποιήσεις που δεν αναπαράγονται πλήρως εδώ \cite{kociembaCubeExplorer}.

\begin{figure}[htbp]\centering
% Ροή του αλγορίθμου δύο φάσεων του Kociemba (περιεχόμενο σχήματος· το
% περιβάλλον figure και η λεζάντα ορίζονται στο κεφάλαιο 4). Φασικά βάθη
% κατά τα κλειδωμένα δεδομένα: φάση 1 <= 12, φάση 2 <= 18.
\begin{tikzpicture}[
  state/.style={draw, rectangle, align=center, font=\small,
                minimum height=1.1cm, minimum width=2.7cm, inner sep=5pt},
  phase/.style={font=\footnotesize},
  desc/.style={font=\scriptsize, align=center, anchor=north},
  flow/.style={-{Stealth[length=2.8mm]}, semithick},
  node distance=2.6cm
]
\node[state] (start) {Αυθαίρετη\\ κατάσταση};
\node[state, right=of start] (g1) {Υποομάδα \(G_1\)};
\node[state, right=of g1] (solved) {Λυμένη\\ κατάσταση};

\draw[flow] (start) -- node[phase, above] {Φάση 1} (g1);
\draw[flow] (g1) -- node[phase, above] {Φάση 2} (solved);

\node[desc] at ($(start.east)!0.5!(g1.west) + (0,-0.7)$)
  {IDA* στις συντεταγμένες\\ φάσης 1, βάθος \(\le 12\)};
\node[desc] at ($(g1.east)!0.5!(solved.west) + (0,-0.7)$)
  {IDA* στις συντεταγμένες\\ φάσης 2, βάθος \(\le 18\)};

\draw[flow, dashed] (solved.north) -- ++(0,0.85) -| (start.north);
\node[font=\scriptsize, align=center, anchor=south]
  at ($(start.north)!0.5!(solved.north) + (0,0.95)$)
  {επαναληπτική εμβάθυνση στο μήκος της φάσης 1\\
   \(\rightarrow\) διαδοχικά καλύτερες συνολικές λύσεις};
\end{tikzpicture}

\caption{Ο αλγόριθμος δύο φάσεων του Kociemba: η φάση 1 (βάθος έως 12) οδηγεί την κατάσταση στην υποομάδα \(G_1\) και η φάση 2 (βάθος έως 18) ολοκληρώνει την επίλυση· η επαναληπτική εμβάθυνση στο μήκος της φάσης 1 παράγει διαδοχικά καλύτερες συνολικές λύσεις.}
\label{fig:kociemba-twophase}
\end{figure}

Η παρούσα εργασία περιλαμβάνει δύο υλοποιήσεις Kociemba. Η πρώτη είναι εγγενής υλοποίηση δύο φάσεων με φραγμένα όρια αναζήτησης. Η φάση 1 χρησιμοποιεί ως παραδεκτό κάτω φράγμα τον ακριβή, συμμετρικά ανηγμένο πίνακα αποστάσεων φάσης 1 βάθους 12. Ελλείψει αυτού επανέρχεται στο μέγιστο των παραγόμενων πινάκων προσανατολισμού γωνιών, προσανατολισμού ακμών και φέτας UD. Επιπλέον συλλέγει περισσότερες από μία ενδιάμεσες καταστάσεις. Οι υποψήφιες καταστάσεις ταξινομούνται με βάση το κάτω φράγμα της φάσης 2, ώστε ο επιλυτής να μη δεσμεύεται στην πρώτη τυχαία μετάβαση προς την υποομάδα. Η φάση 2 εκτελεί αναζήτηση με τα δέκα πραγματικά σύμβολα HTM \(U,U',U2,D,D',D2,R2,L2,F2,B2\), δηλαδή τις δυνάμεις των γεννητόρων \(U,D\) και τις μισές στροφές των \(R,L,F,B\). Ως παραδεκτό φράγμα χρησιμοποιεί πρόσθετες προβολές μετάθεσης γωνιών, μετάθεσης ακμών U/D και μετάθεσης των ακμών της φέτας. Η δεύτερη υλοποίηση είναι προαιρετικός εξωτερικός προσαρμογέας (\eng{adapter}) Kociemba \cite{muodovKociembaPkg}, ο οποίος χρησιμοποιείται μόνο ως πρακτική γραμμή σύγκρισης. Και στις δύο περιπτώσεις κάθε αποτέλεσμα επαληθεύεται από τον ανεξάρτητο επαληθευτή, αλλά ο χαρακτηρισμός παραμένει \code{non_exact}, επειδή δεν αποδεικνύεται καθολική βελτιστότητα.

Η φάση 1 έχει σαφές κατηγόρημα στόχου (\eng{goal predicate}). Η κατάσταση ανήκει στον στόχο όταν οι συντεταγμένες CO, EO και UD-slice είναι ίσες με τις λυμένες τιμές τους. Η IDA* αυξάνει το όριο με βάση την ακριβή, συμμετρικά ανηγμένη απόσταση φάσης 1 ή, ελλείψει του αντίστοιχου πίνακα, τη μέγιστη προβολική απόσταση των τριών αυτών πινάκων. Αν η αναζήτηση βρει ακολουθίες, κάθε ακολουθία είναι έγκυρη μετάβαση προς την ενδιάμεση ομάδα μέσα στα καθορισμένα όρια, χωρίς όμως να αποδεικνύει ότι οδηγεί στη συνολικά μικρότερη λύση.

Η φάση 2 είναι πιο απαιτητική επειδή, παρότι το σύνολο κινήσεων είναι μικρότερο, ο χώρος των μεταθέσεων παραμένει μεγάλος. Η εργασία προσθέτει τρεις φασικές προβολές. Η πρώτη βλέπει τη μετάθεση των γωνιών. Η δεύτερη βλέπει τη μετάθεση των οκτώ ακμών που ανήκουν στις στρώσεις U/D. Η τρίτη βλέπει τη μετάθεση των τεσσάρων ακμών της φέτας. Όλες παράγονται με το ίδιο περιορισμένο σύνολο των δέκα συμβόλων της φάσης 2. Η ευρετική της φάσης 2 παίρνει το μέγιστο από αυτές τις προβολές και από το γενικό φράγμα λανθασμένων κυβιδίων.

Η επιλογή του μέγιστου είναι συντηρητική. Μια ισχυρότερη υλοποίηση θα μπορούσε να χρησιμοποιήσει συνδυασμένες συντεταγμένες ή προσθετικές τεχνικές όπου επιτρέπεται. Η παρούσα εργασία δεν το κάνει στην πρακτική υλοποίηση, επειδή ο στόχος είναι να παραμείνει η παραγωγή πινάκων εφικτή και ελέγξιμη σε συμβατικό υπολογιστή. Το τίμημα είναι ότι η φάση 2 μπορεί να επιστρέψει \code{lower_bound} ή \code{timeout}, όταν τα καθορισμένα όρια χρόνου ή κόμβων είναι στενά για τη συγκεκριμένη κατάσταση. Τα προεπιλεγμένα φασικά βάθη είναι ίσα με τις φασικές διαμέτρους, δηλαδή έως 12 κινήσεις για τη φάση 1 και έως 18 για τη φάση 2. Με αυτά τα όρια ο επιλυτής λύνει τυπικά τυχαία ανακατέματα (\eng{scrambles}), χωρίς αυτό να αποτελεί εγγύηση για κάθε δυνατή κατάσταση. Τα αποτελέσματα αυτά είναι μέρος της καταγραφής και όχι σφάλμα προς απόκρυψη.

Μια ισχυρότερη εκδοχή της ίδιας μηχανικής, προσανατολισμένη στην απόδειξη βελτιστότητας, υλοποιείται χωριστά. Βασίζεται σε συμμετρικά ανηγμένο πίνακα αποστάσεων φάσης 1 και στο παραδεκτό κάτω φράγμα τριών αξόνων του Reid \cite{reid1997optimal}. Η απόσταση φάσης 1 υπολογίζεται ως προς καθέναν από τους τρεις άξονες του κύβου και λαμβάνεται το μέγιστο. Το μέγιστο αυτό παραμένει παραδεκτό, επειδή καμία από τις τρεις προβολές δεν υπερεκτιμά την πραγματική απόσταση. Το φράγμα αυτό τροφοδοτεί εξαντλητική αναζήτηση IDA* ενός ορίου. Η υλοποίηση αυτής της μηχανής περιγράφεται στο Κεφάλαιο~\ref{chap:implementation}, ενώ οι μετρήσεις της, μεταξύ των οποίων η απόδειξη με ίδιο κώδικα ότι η κατάσταση superflip έχει απόσταση ακριβώς 20, παρουσιάζονται στο Κεφάλαιο~\ref{chap:experiments}. Το κάτω φράγμα των 20 κινήσεων για το superflip είχε τεκμηριωθεί ιστορικά από τον Reid \cite{reid1995superflip}.

Ο εξωτερικός προσαρμογέας Kociemba έχει διαφορετικό ρόλο. Δεν τεκμηριώνει την εγγενή υλοποίηση, αλλά προσφέρει πρακτικό σημείο σύγκρισης για το τι μπορεί να επιστρέψει μια ώριμη υλοποίηση δύο φάσεων. Η εργασία τον διατηρεί με χαρακτηρισμό \code{non_exact}, ακόμη και όταν λύνει όλες τις περιπτώσεις του σώματος μετρήσεων, επειδή η τοπική επαλήθευση ελέγχει την εγκυρότητα της λύσης αλλά όχι την καθολική ελαχιστότητα.

\section{Ο αλγόριθμος του Korf και η IDA*}
\label{sec:algorithms-korf}
Ο Korf έδειξε ότι η IDA* με ισχυρούς πίνακες προτύπων (\eng{pattern databases}, PDB) μπορεί να βρει βέλτιστες λύσεις σε τυχαίες περιπτώσεις του κύβου \cite{korf1997pattern}. Η παρούσα υλοποίηση διατηρεί τη δομή της IDA* και χρησιμοποιεί παραδεκτά φράγματα από τέσσερις πηγές. Η πρώτη είναι ένα απλό φράγμα από το πλήθος των λανθασμένων κυβιδίων. Η δεύτερη είναι οι μικροί παραγόμενοι πίνακες προβολών CO/EO/UD-slice. Η τρίτη είναι ένας πίνακας προτύπων για ολόκληρη την κατάσταση των γωνιών του 3×3×3, παραγόμενος σε C/C++. Η τέταρτη είναι οκτώ εξακμικοί πίνακες προτύπων, επίσης παραγόμενοι σε C/C++. Ο γωνιακός πίνακας έχει \(8!\cdot3^7=88\,179\,840\) προβεβλημένες καταστάσεις. Κάθε εξακμικός πίνακας έχει \(\binom{12}{6}\cdot6!\cdot2^6=42\,577\,920\) καταστάσεις, οπότε τα οκτώ αρχεία ακμών καταγράφουν συνολικά \(340\,623\,360\) προβεβλημένες καταστάσεις ακμών. Όλα αποθηκεύονται ως δυαδικοί πίνακες ενός byte ανά κατάσταση. Ο επιλυτής Korf δεν αποτελεί συνεπώς απλή επίδειξη μικρών προβολών σε Python, αλλά πραγματική μηχανή ακριβούς αναζήτησης στον 3×3×3 με πίνακες προτύπων παραγόμενους σε C/C++. Παραμένει όμως ακριβής μόνο όταν η IDA* ολοκληρώνει την αναζήτηση μέσα στα όρια βάθους, χρόνου και κόμβων που έχουν καταγραφεί στο αποτέλεσμα.

Η IDA* ξεκινά από την τιμή της ευρετικής στην αρχική κατάσταση. Για κάθε όριο εκτελεί αναζήτηση κατά βάθος (\eng{depth-first}) και απορρίπτει κλάδους όπου το \(g+h\) ξεπερνά το όριο. Αν δεν βρεθεί λύση, το επόμενο όριο γίνεται η ελάχιστη τιμή που υπερέβη το προηγούμενο όριο. Η διαδικασία αυτή διατηρεί χαμηλές απαιτήσεις μνήμης σε σχέση με τον A*, αλλά επαναλαμβάνει μέρος της εργασίας σε κάθε επανάληψη. Στον κύβο η επανάληψη είναι αποδεκτή μόνο όταν η ευρετική αποκόπτει αρκετούς κλάδους.

Το απλό φράγμα λανθασμένων κυβιδίων βασίζεται στο ότι μία κίνηση HTM επηρεάζει περιορισμένο αριθμό γωνιών και ακμών. Αν υπάρχουν πολλοί λανθασμένοι προσανατολισμοί ή λανθασμένες θέσεις, απαιτούνται τουλάχιστον μερικές κινήσεις. Το φράγμα αποδεικνύεται εύκολα παραδεκτό, αλλά συχνά δίνει μικρές τιμές. Τα φράγματα που βασίζονται σε πίνακες είναι ισχυρότερα, επειδή αποθηκεύουν πραγματικές αποστάσεις προβολών. Ο γωνιακός πίνακας προτύπων κρατά ταυτόχρονα μετάθεση και προσανατολισμό γωνιών, ενώ οι εξακμικοί πίνακες κρατούν θέση, σχετική μετάθεση και προσανατολισμό έξι επιλεγμένων ακμών. Η βασική υλοποίηση παίρνει το μέγιστο των τιμών πλήρους κόστους. Αυτό είναι συντηρητικό και παραδεκτό ακόμη και για επικαλυπτόμενες προβολές ακμών. Η υλοποίηση περιλαμβάνει επίσης χωριστή εκδοχή με κατανομή κόστους (\eng{cost partitioning}): δύο συμβατές προβολές URF/DLB και άθροισμα των αντίστοιχων αποστάσεων με κόστος τελεστή 0/1. Η εκδοχή αυτή είναι παραδεκτή λόγω της κατανομής του κόστους ανά κίνηση HTM. Στο συγκεκριμένο δείγμα, ωστόσο, η αποθηκευμένη μέτρηση κάλυψης δεν έδειξε βελτίωση έναντι του μεγίστου των οκτώ πλήρους κόστους εξακμικών πινάκων.

Το σενάριο \code{scripts/compare_heuristics.py} καταγράφει χωριστά τις τέσσερις συνιστώσες της προεπιλεγμένης ευρετικής και το τελικό μέγιστο που περνά στην IDA*. Στο αποθηκευμένο τεκμήριο της εργασίας, \path{results/processed/heuristic_comparison_seed_2026_thesis.json}, οι ρηχές περιπτώσεις έχουν ανεξάρτητη απόσταση BFS και όλες οι συνιστώσες παραμένουν μικρότερες ή ίσες από αυτήν. Το ίδιο τεκμήριο δείχνει ότι το \code{combined_table_lower_bound} δεν είναι ασθενέστερο από καμία συνιστώσα του. Η τρίτη απαίτηση του θέματος είναι μια κατάλληλη ευρετική συνάρτηση για τον A* ή παραλλαγή του. Έτσι η απαίτηση αυτή συνδέεται με συγκεκριμένο κώδικα και παραγόμενα τεκμήρια, και όχι μόνο με περιγραφική διατύπωση.

Η υλοποίηση επιστρέφει \code{exact} μόνο όταν η IDA* βρει λύση πριν από την εξάντληση του ορίου χρόνου και μέσα στο μέγιστο βάθος. Αν το αρχικό κάτω φράγμα ξεπερνά το επιτρεπτό βάθος, ή αν εξαντληθεί ο χρόνος ή το όριο κόμβων, το αποτέλεσμα είναι \code{timeout}. Ο χαρακτηρισμός \code{timeout} εδώ δεν σημαίνει ότι η κατάσταση δεν λύνεται. Σημαίνει ότι η συγκεκριμένη αναζήτηση, με τα καθορισμένα όρια, δεν παρήγαγε απόδειξη. Αυτή η ακριβής χρήση της ορολογίας είναι κρίσιμη για τη σύνδεση με το θέμα της απόστασης.

\subsection{Προαιρετικοί επταακμικοί πίνακες προτύπων}
\label{sec:algorithms-seven-edge}
Τα προεπιλεγμένα τεκμήρια του επιλυτή Korf στην παρούσα εργασία βασίζονται αποκλειστικά στον γωνιακό πίνακα προτύπων και στους οκτώ εξακμικούς πίνακες, με συνολικό μέγεθος αρχείων 428\,803\,832 bytes. Πέρα από αυτά, η υλοποίηση υποστηρίζει προαιρετικούς επταακμικούς πίνακες προτύπων (\eng{7-edge PDBs}), με \(\binom{12}{7}\cdot7!\cdot2^7=510\,935\,040\) καταστάσεις ανά πίνακα και μέγιστη αποθηκευμένη απόσταση 11. Οι πίνακες αυτοί ενεργοποιούνται ρητά και δεν αποτελούν μέρος των προεπιλεγμένων τεκμηρίων· όταν απουσιάζουν, ο επιλυτής και η ευρετική επανέρχονται ακριβώς στη συμπεριφορά των εξακμικών πινάκων. Η αποθηκευμένη μέτρηση ισχύος, \path{results/processed/seven_edge_strength_seed_2026.json}, δείχνει μέση αύξηση 0.5 κινήσεων του κάτω φράγματος στις τυχαίες βαθιές καταστάσεις του δείγματος, αλλά μηδενική βελτίωση για την κατάσταση superflip, όπου το επταακμικό φράγμα παραμένει 8, ίδιο με το εξακμικό. Ορισμένες καταγεγραμμένες διερευνητικές εκτελέσεις χρησιμοποιούν δέκα πίνακες ακμών (\code{edge_pdb_count=10}, δηλαδή οκτώ εξακμικούς και δύο επταακμικούς)· τα μεταδεδομένα κάθε εκτέλεσης καταγράφουν ρητά τη σύνθεση της ευρετικής, ώστε οι εκτελέσεις αυτές να μη συγχέονται με τα αποτελέσματα που στηρίζονται στην προεπιλεγμένη ευρετική.

\section{Η πλήρης μελέτη του κύβου 2×2×2}
\label{sec:algorithms-pocket}
Η κύρια αλγοριθμική έμφαση της εργασίας παραμένει στον πλήρη κύβο 3×3×3: εγγενής Kociemba, εγγενής Thistlethwaite και Korf/IDA* με γωνιακό πίνακα προτύπων και εξακμικούς πίνακες προτύπων. Η πλήρης μελέτη του 2×2×2 διατηρείται ως μικρό, εξαντλητικό παράδειγμα, για να φανεί πώς μοιάζει μια πλήρης κατανομή αποστάσεων όταν ολόκληρος ο χώρος χωρά τοπικά. Ο κύβος 2×2×2 χρησιμοποιεί μόνο τις οκτώ γωνίες και αποτελεί φυσικό μικρότερο πεδίο για πλήρη αναζήτηση, με ανεξάρτητες τεχνικές περιγραφές σε εκπαιδευτικό υλικό κύβων \cite{randelshoferPocketCube}. Η κατάσταση κανονικοποιείται με σταθερή τη γωνία DBL ως σημείο αναφοράς και χρησιμοποιούνται κινήσεις HTM πάνω στις έδρες \(U\), \(R\) και \(F\). Ο χώρος έχει \(7!\cdot3^6=3\,674\,160\) καταστάσεις. Η BFS πάνω σε αυτόν τον χώρο παράγει πλήρη κατανομή αποστάσεων και αντιπροσωπευτικές βέλτιστες λύσεις. Η μελέτη αποδεικνύει βελτιστότητα μόνο για την κανονικοποιημένη περίπτωση 2×2×2 και δεν μεταφέρεται ως ισχυρισμός πλήρους βελτιστότητας για τον κύβο 3×3×3.

Η κανονικοποίηση με σταθερή γωνία DBL αφαιρεί την ολική περιστροφή του κύβου από τον χώρο καταστάσεων. Οι κινήσεις \(U,R,F\) αρκούν ως γεννήτριες για το κανονικοποιημένο μοντέλο. Η κατάσταση κωδικοποιείται ως μετάθεση επτά γωνιών και έξι ανεξάρτητοι προσανατολισμοί. Ο πλήρης πίνακας αποστάσεων αποθηκεύεται κατά την BFS ως πίνακας μικρών ακεραίων, ώστε να είναι εφικτή η απαρίθμηση χωρίς να αποθηκεύονται γονείς για κάθε κατάσταση.

Για την πλήρη κατανομή δεν χρειάζεται να ανακτηθεί διαδρομή προς κάθε κατάσταση· αρκεί η απόσταση. Για τις αντιπροσωπευτικές λύσεις, όμως, εκτελείται BFS με καταγραφή γονέων ή διαδρομής, ώστε να παραχθεί συγκεκριμένη ακολουθία. Ο διαχωρισμός αυτός μειώνει τη μνήμη της πλήρους απαρίθμησης και διατηρεί την αναγνωσιμότητα των παραδειγμάτων. Οι αντιπροσωπευτικές γραμμές αποτελεσμάτων επαληθεύονται χωριστά από το σενάριο επαλήθευσης αποτελεσμάτων.

Η μελέτη του 2×2×2 έχει διπλό ρόλο. Πρώτον, προσφέρει μέσα στην εργασία μια πλήρη περίπτωση αποδεδειγμένης βελτιστότητας. Δεύτερον, λειτουργεί ως μικρογραφία της μεθοδολογίας που θα απαιτούνταν για ισχυρότερη απόδειξη στον 3×3×3: σαφής χώρος, πλήρες σύνολο γεννητριών κινήσεων, αποθηκευμένες αποστάσεις, έλεγχοι από τον επαληθευτή και αυτόματα παραγόμενοι πίνακες LaTeX. Δεν αντικαθιστά τους αλγορίθμους του 3×3×3, αλλά συμπληρώνει την εργασία εκεί όπου η πλήρης βελτιστότητα στον 3×3×3 παραμένει εκτός τοπικού ορίου.

\section{Αναγνώριση απόστασης}
\label{sec:algorithms-distance}
Η εντολή \code{distance} επιστρέφει μία από τις κατηγορίες \code{exact_distance}, \code{lower_bound}, \code{unknown_timeout} ή \code{invalid_state}. Στα πειράματα του Κεφαλαίου~\ref{chap:experiments} οι ρηχές περιπτώσεις αποδεικνύονται με BFS, ενώ η τυχαία περίπτωση του γρήγορου σώματος μετρήσεων επιστρέφει μόνο κάτω φράγμα, όταν η ακριβής αναζήτηση δεν ολοκληρώνεται. Η ίδια εντολή μπορεί επίσης να ενεργοποιήσει τη μηχανή H48, ώστε μια κατάσταση σε μορφή ψηφίδων (\eng{facelets}) ή κυβιδίων να δώσει άμεσα ακριβή απόσταση, όταν το υποσύστημα H48 ολοκληρώσει την αναζήτηση.

Η σειρά ενεργειών είναι συντηρητική. Πρώτα ελέγχεται η φυσική εγκυρότητα. Αν η κατάσταση είναι άκυρη, η απόκριση είναι \code{invalid_state} και δεν εκτελείται αναζήτηση. Έπειτα δοκιμάζεται BFS μέχρι καθορισμένο βάθος. Αν βρεθεί λύση, η απόσταση είναι ακριβής, επειδή έχουν εξεταστεί όλα τα μικρότερα βάθη. Αν ζητηθεί ο εγγενής βέλτιστος επιλυτής ή η μηχανή H48, η εντολή χρησιμοποιεί την αντίστοιχη μηχανή ακριβούς αναζήτησης μετά την BFS. Μόνο αν αυτή ολοκληρώσει λύση και η λύση επαληθευτεί, επιστρέφεται \code{exact_distance}. Διαφορετικά επιστρέφεται κάτω φράγμα ή \code{unknown_timeout}.

Η αναγνώριση απόστασης δεν πρέπει να συγχέεται με την πρακτική επίλυση. Ένας επιλυτής τύπου Kociemba μπορεί να επιστρέψει γρήγορα λύση χωρίς να αποδείξει το ελάχιστο μήκος. Η εντολή \code{distance} έχει διαφορετική ευθύνη: να δηλώσει τι γνωρίζει με απόδειξη για την απόσταση. Γι' αυτό μπορεί να είναι πιο αυστηρή και να επιστρέψει μόνο κάτω φράγμα εκεί όπου ένας πρακτικός επιλυτής βρίσκει λύση. Η αυστηρότητα αυτή είναι αναγκαία λόγω της διατύπωσης του θέματος.

\section{Συνδυασμός επιλυτών στις μετρήσεις}
\label{sec:algorithms-solver-mix}
Το κεντρικό σώμα μετρήσεων της εργασίας δεν αναζητά έναν μοναδικό νικητή. Χρησιμοποιεί πολλούς επιλυτές για να αναδείξει διαφορετικές εγγυήσεις. Ο επιλυτής Korf/IDA* αντιπροσωπεύει τη βέλτιστη αναζήτηση με παραδεκτό φράγμα σε εφικτά βάθη. Ο εγγενής Kociemba αντιπροσωπεύει πρακτική επίλυση διαδοχικών σταδίων με πίνακες φάσεων. Ο προσαρμογέας αντιπροσωπεύει ώριμη εξωτερική πρακτική υλοποίηση δύο φάσεων. Ο εγγενής Thistlethwaite αντιπροσωπεύει τετραφασική εκπαιδευτική αλυσίδα υποομάδων με τελικό στάδιο μισών στροφών. Κάθε γραμμή αποτελεσμάτων καταγράφει την ίδια αρχική κατάσταση, το ίδιο βάθος ανακατέματος και το ίδιο συμβόλαιο επαλήθευσης.

Η σχεδίαση αυτή επιτρέπει να διαβαστούν τα αποτελέσματα χωρίς υπεραπλούστευση. Αν ο προσαρμογέας ή οι εγγενείς επιλυτές διαδοχικών σταδίων λύνουν όλες τις περιπτώσεις, αυτό δείχνει πρακτική ισχύ στο συγκεκριμένο σώμα μετρήσεων· δεν δείχνει βέλτιστη απόσταση. Αν ο επιλυτής Korf καταγράφει αποτελέσματα \code{timeout}, αυτό δείχνει το κόστος της απόδειξης· δεν σημαίνει ότι η κατάσταση είναι δύσκολη για κάθε πρακτικό επιλυτή. Η σύγκριση επομένως διαχωρίζει την πρακτική επίλυση από την ακριβή απόδειξη απόστασης.
